\documentclass[aspectratio=169,10pt]{beamer}
\makeatletter\def\input@path{{../estilo/}{estilo/}}\makeatother
\usepackage{devmobile}
\modulo{12}{Publicação na Google Play}{Build de release, assinatura, App Bundle, Play Console, faixas de teste, verificação de desenvolvedor, LGPD e CI/CD}

\begin{document}
\capa

\begin{frame}{Objetivos do módulo}
  \begin{conceito}[Ao final deste módulo você será capaz de]
    \begin{itemize}
      \item Preparar um build de release otimizado e versionado.
      \item Assinar o app e entender a Play App Signing.
      \item Criar a conta e a ficha do app na Play Console, incluindo as declarações obrigatórias.
      \item Conduzir testes interno, fechado e aberto e cumprir a regra de 12 testadores.
      \item Lançar de forma gradual e acompanhar a saúde do app em produção.
      \item Atender à LGPD e automatizar a publicação com CI/CD.
    \end{itemize}
  \end{conceito}
\end{frame}

\sumario

% =============================================================================
\section{Preparando o release}

\begin{frame}[fragile]{Build de release}
  \begin{columns}[T]
    \begin{column}{0.54\textwidth}
\begin{codigo}[title=app/build.gradle.kts]{Kotlin}
android {
  defaultConfig {
    versionCode = 7          // inteiro, sempre crescente
    versionName = "1.2.0"    // exibido ao usuário
  }
  buildTypes {
    release {
      isMinifyEnabled = true      // R8: encolhe e ofusca
      isShrinkResources = true    // remove recursos não usados
      proguardFiles(getDefaultProguardFile(
          "proguard-android-optimize.txt"),
        "proguard-rules.pro")
      signingConfig = signingConfigs.getByName("release")
    }
  }
}
\end{codigo}
    \end{column}
    \begin{column}{0.44\textwidth}
      \begin{conceito}[Versionamento]
        \cod{versionCode} identifica a versão para a loja: cada envio precisa de um número maior. \cod{versionName} segue o \emph{versionamento semântico} (maior.menor.correção).
      \end{conceito}
      \begin{atencao}[R8]
        A ofuscação renomeia classes. Bibliotecas que usam reflexão podem precisar de regras \cod{-keep}. \textbf{Teste o build release} antes de enviar.
      \end{atencao}
      \begin{dica}
        Guarde o \arq{mapping.txt} de cada versão: ele desofusca as pilhas de erro no Crashlytics e na Play Console.
      \end{dica}
    \end{column}
  \end{columns}
\end{frame}

\begin{frame}{Checklist antes de publicar}
  \begin{columns}[T]
    \begin{column}{0.48\textwidth}
      \begin{conceito}[Produto]
        \begin{itemize}
          \item Ícone adaptativo (\emph{Image Asset Studio}) e nome definitivos.
          \item \cod{applicationId} definitivo (não muda depois).
          \item \cod{targetSdk} atendendo à exigência vigente.
          \item Textos em \arq{strings.xml}; tema claro e escuro revisados.
          \item Mensagens de erro, estados vazios e modo offline testados.
        \end{itemize}
      \end{conceito}
    \end{column}
    \begin{column}{0.48\textwidth}
      \begin{conceito}[Técnico]
        \begin{itemize}
          \item Build release testado em aparelhos reais (versões antigas e novas).
          \item Logs de depuração e URLs de teste removidos.
          \item Nenhum segredo no código ou no repositório.
          \item Crashlytics ativo e \arq{mapping.txt} enviado.
          \item Apenas as permissões realmente usadas.
        \end{itemize}
      \end{conceito}
    \end{column}
  \end{columns}
  \begin{atencao}[Nível de API alvo]
    Desde \textbf{31/08/2026}, apps novos e atualizações na Google Play precisam de \cod{targetSdk} \textbf{36} (Android 16) ou superior. A exigência sobe todo ano, por volta de agosto.
  \end{atencao}
\end{frame}

% =============================================================================
\section{Assinatura do app}

\begin{frame}{Play App Signing}
  \begin{center}
  \begin{tikzpicture}[node distance=9mm]
    \node[caixaescura,minimum width=2.8cm] (dev) {Você\\ \tiny gera o AAB};
    \node[caixaambar,minimum width=2.8cm,right=of dev] (up) {AAB assinado com a\\ \textbf{chave de upload}};
    \node[caixapetroleo,minimum width=2.8cm,right=of up] (gp) {Google Play\\ \tiny verifica e reassina};
    \node[caixa,minimum width=2.8cm,right=of gp] (apk) {APKs assinados com a\\ \textbf{chave do app}};
    \draw[seta] (dev) -- (up); \draw[seta] (up) -- node[rotulo,above]{envio} (gp); \draw[seta] (gp) -- node[rotulo,above]{instalação} (apk);
  \end{tikzpicture}
  \end{center}
  \begin{columns}[T]
    \begin{column}{0.48\textwidth}
      \begin{conceito}[Duas chaves]
        A \textbf{chave do app} fica guardada pelo Google e assina o que chega aos usuários. A \textbf{chave de upload} fica com você e só prova ao Google que o envio é seu.
      \end{conceito}
    \end{column}
    \begin{column}{0.48\textwidth}
      \begin{dica}[Vantagem]
        Se a chave de upload vazar ou for perdida, você pede a troca ao suporte da Play. Antes da Play App Signing, perder a chave significava nunca mais atualizar o app.
      \end{dica}
    \end{column}
  \end{columns}
  \begin{atencao}
    O SHA-1 da chave do app (na Play Console) precisa ser cadastrado no Firebase para o login Google funcionar em produção (Módulo 09).
  \end{atencao}
\end{frame}

\begin{frame}[fragile]{Criando a chave e assinando}
  \begin{columns}[T]
    \begin{column}{0.54\textwidth}
\begin{terminal}[title=Gerando a chave de upload]
keytool -genkeypair -v -keystore upload.jks \
  -alias upload -keyalg RSA -keysize 2048 \
  -validity 10000
\end{terminal}
\begin{codigo}[title=app/build.gradle.kts]{Kotlin}
signingConfigs {
  create("release") {
    storeFile = file(System.getenv("KEYSTORE_PATH")
      ?: "upload.jks")
    storePassword = System.getenv("KEYSTORE_PASS")
    keyAlias = "upload"
    keyPassword = System.getenv("KEY_PASS")
  }
}
\end{codigo}
\begin{terminal}[title=Gerando o App Bundle]
./gradlew bundleRelease
# app/build/outputs/bundle/release/app-release.aab
\end{terminal}
    \end{column}
    \begin{column}{0.44\textwidth}
      \begin{erro}[Nunca no Git]
        O \arq{.jks} e as senhas \textbf{não} vão para o repositório. Use variáveis de ambiente, \arq{\textasciitilde/.gradle/gradle.properties} ou \emph{secrets} da CI.
      \end{erro}
      \begin{dica}[Backup]
        Guarde o \arq{.jks} e as senhas num gerenciador de senhas e em uma cópia offline.
      \end{dica}
      \begin{dica}
        Pelo Android Studio: \textbf{Build $\rightarrow$ Generate Signed App Bundle / APK} faz o mesmo com um assistente.
      \end{dica}
    \end{column}
  \end{columns}
\end{frame}

% =============================================================================
\section{Play Console}

\begin{frame}{A conta de desenvolvedor}
  \begin{columns}[T]
    \begin{column}{0.5\textwidth}
      \scriptsize
      \begin{tabular}{@{}p{1.9cm} p{2.0cm} p{2.3cm}@{}}
        \toprule
        & \cabecalho{Pessoal} & \cabecalho{Organização} \\
        \midrule
        Para & pessoa física & empresas, instituições \\
        Exige & verificação de identidade & D-U-N-S e verificação \\
        Teste fechado obrigatório & sim (contas novas) & não \\
        Nome exibido & seu nome & nome da organização \\
        \bottomrule
      \end{tabular}
      \vspace{3pt}
      \begin{conceito}[Custo]
        Taxa de registro \textbf{única}, paga ao criar a conta em \url{https://play.google.com/console}.
      \end{conceito}
    \end{column}
    \begin{column}{0.48\textwidth}
      \begin{dica}[Para projetos institucionais]
        Apps que representam o IFMA, um laboratório ou um projeto de extensão ficam melhor em uma \textbf{conta de organização} da instituição, que garante continuidade quando os alunos se formam.
      \end{dica}
      \begin{atencao}
        A conta exige verificação de identidade, telefone e e-mail de contato públicos. Mantenha esses dados atualizados: contas inativas ou com dados inválidos podem ser encerradas.
      \end{atencao}
    \end{column}
  \end{columns}
\end{frame}

\begin{frame}{Ficha da loja e declarações obrigatórias}
  \begin{columns}[T]
    \begin{column}{0.48\textwidth}
      \begin{conceito}[Ficha da loja]
        \begin{itemize}
          \item \textbf{Nome} do app (até 30 caracteres).
          \item \textbf{Descrição curta} (até 80) e \textbf{completa} (até 4000).
          \item \textbf{Ícone} 512$\times$512\,px.
          \item \textbf{Gráfico de recursos} 1024$\times$500\,px.
          \item \textbf{Capturas de tela} (celular e tablet).
          \item Categoria e e-mail de contato.
        \end{itemize}
      \end{conceito}
    \end{column}
    \begin{column}{0.48\textwidth}
      \begin{conceito}[Conteúdo do app (Política)]
        \begin{itemize}
          \item URL da \textbf{política de privacidade}.
          \item Formulário de \textbf{segurança dos dados}.
          \item \textbf{Classificação de conteúdo} (questionário IARC).
          \item \textbf{Público-alvo} e se contém anúncios.
          \item Declaração de permissões sensíveis.
          \item Acesso para o revisor (se houver login).
        \end{itemize}
      \end{conceito}
    \end{column}
  \end{columns}
  \begin{atencao}
    Se o app permite criar conta, ele deve oferecer \textbf{exclusão da conta} dentro do app e por um link na web.
  \end{atencao}
\end{frame}

% =============================================================================
\section{Faixas de teste}

\begin{frame}{Do teste interno à produção}
  \begin{center}
  \begin{tikzpicture}[node distance=6mm]
    \node[caixa,minimum width=2.6cm,minimum height=12mm] (i) {\textbf{Interno}\\ \tiny até 100 testadores\\ disponível em minutos};
    \node[caixaambar,minimum width=2.6cm,minimum height=12mm,right=of i] (f) {\textbf{Fechado}\\ \tiny lista de e-mails ou\\ Grupo Google};
    \node[caixapetroleo,minimum width=2.6cm,minimum height=12mm,right=of f] (a) {\textbf{Aberto}\\ \tiny qualquer pessoa pode\\ entrar no teste};
    \node[caixaescura,minimum width=2.6cm,minimum height=12mm,right=of a] (p) {\textbf{Produção}\\ \tiny todos os usuários};
    \draw[seta] (i) -- (f); \draw[seta] (f) -- (a); \draw[seta] (a) -- (p);
    \draw[setaambar,rounded corners] (f.south) -- ++(0,-5mm) -| node[rotulo,pos=0.25,below]{contas pessoais novas: 12 testadores por 14 dias antes de pedir acesso à produção} (p.south);
  \end{tikzpicture}
  \end{center}
  \begin{columns}[T]
    \begin{column}{0.48\textwidth}
      \begin{atencao}[Regra dos 12 testadores]
        Contas \textbf{pessoais} criadas após 13/11/2023 precisam manter um teste fechado com \textbf{pelo menos 12 testadores inscritos por 14 dias seguidos} antes de solicitar acesso à produção.
      \end{atencao}
    \end{column}
    \begin{column}{0.48\textwidth}
      \begin{dica}[Na prática]
        \begin{itemize}
          \item Recrute 15--20 testadores: se alguém sair, a contagem pode recomeçar.
          \item Uma turma do curso é um ótimo grupo de testadores!
          \item Publique atualizações durante o teste e registre o feedback.
        \end{itemize}
      \end{dica}
    \end{column}
  \end{columns}
\end{frame}

% =============================================================================
\section{Lançamento e operação}

\begin{frame}{Lançamento gradual e saúde do app}
  \begin{columns}[T]
    \begin{column}{0.48\textwidth}
      \begin{conceito}[Lançamento gradual]
        Libere a nova versão para 5\,\%, 20\,\%, 50\,\% e depois 100\,\% dos usuários, acompanhando as falhas a cada etapa. Se algo der errado, \textbf{interrompa} o lançamento.
      \end{conceito}
      \begin{conceito}[Android vitals]
        A Play Console mostra a taxa de falhas e de ANRs percebidos pelos usuários. Acima dos limites de ``mau comportamento'', o app perde visibilidade na loja.
      \end{conceito}
    \end{column}
    \begin{column}{0.48\textwidth}
      \begin{dica}[Depois de publicado]
        \begin{itemize}
          \item Responda às avaliações: a resposta é pública.
          \item Use \textbf{In-App Updates} para pedir a atualização de versões críticas.
          \item Use \textbf{In-App Review} para pedir avaliação no momento certo.
          \item Leia as notas de política da Play: as regras mudam.
          \item Planeje a atualização anual do \cod{targetSdk}.
        \end{itemize}
      \end{dica}
    \end{column}
  \end{columns}
\end{frame}

\begin{frame}{Distribuição fora da Play e verificação de desenvolvedor}
  \begin{columns}[T]
    \begin{column}{0.52\textwidth}
      \begin{atencao}[Novidade de 2026 no Brasil]
        Desde \textbf{30/09/2026}, aparelhos Android certificados no Brasil (e também na Indonésia, em Singapura e na Tailândia) só instalam normalmente apps \textbf{registrados por desenvolvedores verificados}, inclusive quando distribuídos fora da Google Play (APK em site, lojas alternativas). A expansão global está prevista para 2027.
      \end{atencao}
      \begin{conceito}[Como se registrar]
        Pela \textbf{Play Console} (apps na loja) ou pela \textbf{Android Developer Console} (apps só fora da loja), verificando a identidade e registrando o nome do pacote e a chave de assinatura.
      \end{conceito}
    \end{column}
    \begin{column}{0.46\textwidth}
      \begin{dica}[Estudantes e hobbistas]
        Existem contas de \textbf{distribuição limitada}, pensadas para estudantes e projetos pessoais que instalam apps em poucos aparelhos.
      \end{dica}
      \begin{conceito}[O que não muda para o desenvolvimento]
        Instalar pelo Android Studio ou via \textbf{ADB} continua funcionando. Usuários avançados também têm um fluxo especial para instalar apps não verificados.
      \end{conceito}
      \begin{dica}
        Para distribuir o LabMonitor só no laboratório via APK, registre o app antes.
      \end{dica}
    \end{column}
  \end{columns}
\end{frame}

% =============================================================================
\section{LGPD e privacidade}

\begin{frame}{LGPD no desenvolvimento de apps}
  \begin{columns}[T]
    \begin{column}{0.5\textwidth}
      \begin{conceito}[Lei Geral de Proteção de Dados (Lei 13.709/2018)]
        Regula o tratamento de \textbf{dados pessoais} no Brasil: qualquer informação que identifique uma pessoa (nome, e-mail, matrícula, localização, identificadores do aparelho).
      \end{conceito}
      \begin{conceito}[Princípios-chave]
        \textbf{Finalidade} (para quê), \textbf{necessidade} (o mínimo), \textbf{transparência} (explicar), \textbf{segurança} (proteger) e uma \textbf{base legal} (ex.: consentimento ou execução de contrato).
      \end{conceito}
    \end{column}
    \begin{column}{0.48\textwidth}
      \begin{dica}[Checklist para o app]
        \begin{itemize}
          \item Colete só o necessário (\emph{privacy by design}).
          \item Política de privacidade clara, no app e na loja.
          \item Consentimento para usos opcionais (analytics).
          \item Acesso, correção e exclusão dos dados e da conta.
          \item HTTPS, criptografia e regras no backend.
        \end{itemize}
      \end{dica}
      \begin{atencao}
        Não é assessoria jurídica: em projetos institucionais, consulte o encarregado de dados (DPO).
      \end{atencao}
    \end{column}
  \end{columns}
\end{frame}

% =============================================================================
\section{Automatizando com CI/CD}

\begin{frame}[fragile]{Publicação automática}
  \begin{columns}[T]
    \begin{column}{0.56\textwidth}
\begin{codigo}[title=.github/workflows/release.yml (trecho)]{bash}
on:
  push:
    tags: ["v*"]          # ex.: git tag v1.2.0
jobs:
  release:
    runs-on: ubuntu-latest
    steps:
      - uses: actions/checkout@v4
      - uses: actions/setup-java@v4
        with: { distribution: temurin, java-version: 17 }
      - name: Restaurar keystore
        run: echo "$KEYSTORE_B64" | base64 -d > upload.jks
        env: { KEYSTORE_B64: "${{ secrets.KEYSTORE_B64 }}" }
      - name: Gerar AAB
        run: ./gradlew bundleRelease
        env:
          KEYSTORE_PASS: ${{ secrets.KEYSTORE_PASS }}
          KEY_PASS: ${{ secrets.KEY_PASS }}
      # enviar à faixa interna (Gradle Play Publisher
      # ou Fastlane supply) com uma conta de serviço
\end{codigo}
    \end{column}
    \begin{column}{0.42\textwidth}
      \begin{conceito}[Pipeline típico]
        \emph{push} de uma tag $\rightarrow$ testes $\rightarrow$ build assinado $\rightarrow$ envio à faixa interna $\rightarrow$ promoção manual para produção.
      \end{conceito}
      \begin{dica}[Ferramentas]
        \textbf{Gradle Play Publisher} (plugin Gradle) ou \textbf{Fastlane} (\cod{supply}) usam a API da Play Console com uma \emph{conta de serviço} do Google Cloud.
      \end{dica}
      \begin{atencao}
        Segredos ficam nos \emph{secrets} do repositório, nunca no YAML.
      \end{atencao}
    \end{column}
  \end{columns}
\end{frame}

% =============================================================================
\section{Fechamento}

\begin{frame}{Resumo}
  \begin{resumo}
    \begin{itemize}
      \item Release = R8, \cod{versionCode} crescente, \cod{targetSdk} atualizado e testes no build final.
      \item Play App Signing: você guarda a chave de upload; o Google guarda a chave do app.
      \item A loja exige ficha completa, política de privacidade, segurança dos dados e classificação.
      \item Faixas interna, fechada e aberta antes da produção; contas pessoais novas precisam de 12 testadores por 14 dias.
      \item No Brasil, desde 30/09/2026, apps precisam ser de desenvolvedores verificados mesmo fora da Play.
      \item LGPD: mínimo de dados, transparência, consentimento e exclusão. CI/CD automatiza tudo.
    \end{itemize}
  \end{resumo}
\end{frame}

\begin{frame}{Exercícios}
  \begin{exercicio}[Exercício 1 --- Release]
    Gere o AAB assinado do LabMonitor, instale o build release com \cod{bundletool} e verifique se tudo funciona com o R8 ativo.
  \end{exercicio}
  \begin{exercicio}[Exercício 2 --- Política]
    Escreva a política de privacidade do LabMonitor e preencha um rascunho do formulário de segurança dos dados com base no que o app realmente coleta.
  \end{exercicio}
  \begin{exercicio}[Exercício 3 --- Ficha]
    Produza ícone, gráfico de recursos, 4 capturas de tela e as descrições curta e completa do app.
  \end{exercicio}
  \begin{projeto}[LabMonitor --- Etapa 12]
    Publique o app na faixa de \textbf{teste fechado} com a turma como testadores, colete o feedback por 14 dias e lance a versão 1.0.
  \end{projeto}
\end{frame}

\begin{frame}{Para aprofundar}
  \small
  \begin{itemize}
    \item Preparar o app para lançamento: \url{https://developer.android.com/studio/publish/preparing}
    \item Assinatura de apps: \url{https://developer.android.com/studio/publish/app-signing}
    \item Requisitos de nível de API: \url{https://developer.android.com/google/play/requirements/target-sdk}
    \item Testes para contas pessoais novas: \url{https://support.google.com/googleplay/android-developer/answer/14151465}
    \item Verificação de desenvolvedor: \url{https://developer.android.com/developer-verification}
    \item LGPD (texto da lei): \url{https://www.planalto.gov.br/ccivil_03/_ato2015-2018/2018/lei/l13709.htm}
  \end{itemize}
\end{frame}

\encerramento{Módulo 13 --- Desenvolvimento multiplataforma: Flutter, React Native e Kotlin Multiplatform.}

\end{document}
